\section{Problemas Encontrados y Soluciones}

Durante la fase de implementación y pruebas se presentaron diversos problemas que
requerían diagnóstico y corrección. Esta sección documenta cada incidente, su
causa raíz, la solución aplicada y el estado final.

\subsection{Problemas de Configuración}

\subsubsection{Problema 1: SSH No Responde en Routers}

\textbf{Descripción:}

Después de la configuración inicial, se intentó conectar vía SSH al Router 1
desde un cliente externo, pero la conexión fue rechazada.

\textbf{Síntomas:}

\begin{verbatim}
C:\> ssh admin@172.16.192.1

ssh: connect to host 172.16.192.1 port 22: Connection refused
\end{verbatim}

\textbf{Diagnóstico:}

Se accedió al router mediante consola local y se ejecutó:

\begin{verbatim}
Router1# show run | include ssh
ip ssh version 2
ip ssh time-out 120
line vty 0 4
\end{verbatim}

La causa era la seguridad usada en los routers Cisco era deficiente con respecto a lo que Windows 10 soporta, por lo que se necesitaba configurar un cifrado más fuerte o iniciar ssh con un programa como PuTTy.

\subsubsection{Problema 2: DHCP No Asigna Direcciones en LAN 2}

\textbf{Descripción:}

Los clientes en LAN 2 no adquieren direcciones DHCP automáticamente. Se quedan
con 169.254.x.x (direcciones APIPA de respaldo).

\textbf{Síntomas:}

\begin{verbatim}
C:\> ipconfig

Ethernet adapter Ethernet:
   IPv4 Address                     : 169.254.100.50    <--- APIPA, no DHCP
   Subnet Mask                      : 255.255.0.0
\end{verbatim}

\textbf{Diagnóstico:}

Se revisó la configuración del router en LAN 2:

\begin{verbatim}
Router1# show run interface GigabitEthernet0/1

interface GigabitEthernet0/1
 ip address 172.16.196.1 255.255.252.0
 ip dhcp excluded-address 172.16.196.1 172.16.196.10
 ip dhcp pool LAN2
  network 172.16.196.0 255.255.252.0
  default-router 172.16.196.1
  dns-server 172.16.202.1
\end{verbatim}

El pool estaba configurado correctamente. Se revisó la conectividad física:
los switches reportaban que el puerto estaba down.

\textbf{Causa Raíz:}

El cable Ethernet entre el router y el switch de LAN 2 estaba
defectuoso. El puerto del router mostraba \texttt{administratively down}.

\textbf{Solución Aplicada:}

1. Reemplazar el cable Ethernet
2. Verificar que el puerto estuviera activo:

\begin{verbatim}
Router1(config)# interface GigabitEthernet0/1
Router1(config-if)# no shutdown
Router1(config-if)# end
\end{verbatim}

\textbf{Validación Posterior:}

\begin{verbatim}
C:\> ipconfig /renew

IPv4 Address                     : 172.16.196.12
Subnet Mask                      : 255.255.252.0
Default Gateway                  : 172.16.196.1
DHCP Server                      : 172.16.196.1
\end{verbatim}


\subsection{Problemas de Seguridad}

\subsubsection{Problema 6: Telnet Accesible en Router}

\textbf{Descripción:}

Aunque SSH se habilitó, Telnet seguía siendo accesible en el router,
violando los requerimientos de seguridad.

\textbf{Síntomas:}

\begin{verbatim}
C:\> telnet 172.16.192.1

Connected to 172.16.192.1.
User Access Verification
Password: ****   <--- Telnet aún funciona
\end{verbatim}

\textbf{Causa Raíz:}

Se había agregado SSH al transporte, pero no se removió Telnet:

\begin{verbatim}
line vty 0 4
 transport input ssh telnet   <--- Ambos habilitados
\end{verbatim}

\textbf{Solución Aplicada:}

\begin{verbatim}
Router1# configure terminal
Router1(config)# line vty 0 4
Router1(config-line)# transport input ssh
Router1(config-line)# no transport input telnet
Router1(config-line)# end
\end{verbatim}

\textbf{Validación Posterior:}

\begin{verbatim}
C:\> telnet 172.16.192.1

Connecting To 172.16.192.1...Could not open connection to the host,
on port 23: Connect failed

C:\> ssh admin@172.16.192.1

admin@172.16.192.1's password: ****
Router1>
\end{verbatim}

